\documentclass[9pt]{ctexbeamer}

% ====== 清空默认 ======
\setbeamertemplate{navigation symbols}{}
\setbeamertemplate{headline}{}
\setbeamertemplate{footline}{}

% ====== 配色 ======
\definecolor{ink}{HTML}{1B2A4A}
\definecolor{cold}{HTML}{F4F8FF}
\definecolor{ember}{HTML}{E05C2A}
\definecolor{slate}{HTML}{4B5E7D}
\definecolor{rule}{HTML}{BCC9DB}
\definecolor{accent}{HTML}{2E86DE}
\definecolor{soft}{HTML}{EBF1FA}
\definecolor{darkbg}{HTML}{0F1B33}

% ====== 顶部深色条 + 标题 ======
\setbeamercolor{section in head/foot}{fg=white,bg=ink}
\setbeamertemplate{frametitle}{
  \nointerlineskip
  \begin{beamercolorbox}[wd=\paperwidth,ht=0.85cm,dp=0.3cm,leftskip=0.5cm]{section in head/foot}
    {\large\bfseries\insertframetitle}
  \end{beamercolorbox}
}

% ====== 页脚 ======
\setbeamertemplate{footline}{
  \begin{beamercolorbox}[wd=\paperwidth,ht=0.35cm,dp=0.15cm,rightskip=0.4cm]{section in head/foot}
    {\tiny\color{rule}\insertframenumber\,/\,\inserttotalframenumber}
  \end{beamercolorbox}
}

% ====== 内容页背景 ======
\setbeamercolor{background canvas}{bg=cold}

% ====== 列表 ======
\setbeamertemplate{itemize item}{\color{ember}\rule{3pt}{3pt}}
\setbeamertemplate{itemize subitem}{\color{slate}\rule{2pt}{2pt}}

% ====== 自定义 leftblock ======
\newenvironment{leftblock}[1]{%
  \vspace{0.3cm}
  \begin{beamercolorbox}[wd=\textwidth,leftskip=0.4cm,rightskip=0.4cm]{}
    \textcolor{ember}{\rule{2.5pt}{0.55cm}}\hspace{0.2cm}
    {\bfseries\color{ink}#1}
    \par\vspace{0.1cm}
    \hspace{0.35cm}\begin{minipage}{\dimexpr\textwidth-1.2cm\relax}
}{%
    \end{minipage}
  \end{beamercolorbox}
  \vspace{0.15cm}
}

% ====== 包 ======
\usepackage{booktabs}
\usepackage{amsmath}
\usepackage{xcolor}
\usepackage{colortbl}
\usepackage{tikz}
\usepackage{pgfplots}
\pgfplotsset{compat=1.18}
\usetikzlibrary{backgrounds,calc}

% ====== pgfplots 全局 ======
\pgfplotsset{
  every axis/.append style={
    font=\tiny,
    tick label style={font=\fontsize{4.5}{5}\selectfont},
    label style={font=\tiny\color{slate}},
    grid style={dashed,gray!25,line width=0.3pt},
    axis line style={rule,line width=0.4pt},
    legend style={
      font=\fontsize{4.5}{5}\selectfont,
      legend columns=1,
      /tikz/every even column/.append style={column sep=3pt},
      draw=none,
      fill=white,fill opacity=0.7,
      inner sep=1.5pt,
      /tikz/line width=0.6pt,
    },
  },
  barplot/.style={
    ybar, width=0.93\textwidth, height=3.6cm,
    grid=major, bar width=5pt, enlarge x limits=0.12,
  },
  lineplot/.style={
    width=0.93\textwidth, height=3.2cm,
    grid=major, mark=*, mark size=1.2pt, line width=0.8pt,
  },
  narrowbar/.style={
    ybar, width=\columnwidth, height=3.0cm,
    grid=major, bar width=4pt, enlarge x limits=0.18,
  },
  narrowline/.style={
    width=\columnwidth, height=2.8cm,
    grid=major, mark=*, mark size=1pt, line width=0.7pt,
  },
}

% ====== 快捷命令 ======
\newcommand{\tb}{\textcolor{ember}{\textbf{【待补充】}}}
\newcommand{\hi}[1]{\textcolor{ember}{\textbf{#1}}}
\newcommand{\blu}[1]{\textcolor{accent}{\textbf{#1}}}

\begin{document}

% ====== Slide 1: 封面 ======
{
\setbeamercolor{background canvas}{bg=darkbg}
\setbeamertemplate{footline}{}
\setbeamertemplate{frametitle}{}

\begin{frame}[plain]

\begin{tikzpicture}[remember picture,overlay]
  \shade[top color=darkbg,bottom color=ink] (current page.south west) rectangle (current page.north east);
  \foreach \y in {0,2,...,16} {
    \foreach \x in {0,2,...,12} {
      \pgfmathsetmacro\r{random()}
      \pgfmathsetmacro\c{int(\r < 0.42 ? 1 : 0)}
      \node[text=white,opacity=0.09,font=\ttfamily\tiny] at ($(current page.north west)+(1.2+0.85*\x,-0.9-0.6*\y)$) {\c};
    }
  }
  \draw[ember,line width=0.6pt,opacity=0.5] (current page.north west) -- ($(current page.north west)!0.4!(current page.south east)$);
  \draw[accent,line width=0.4pt,opacity=0.35] ($(current page.north east)!0.15!(current page.south east)$) -- ($(current page.south east)!0.3!(current page.south west)$);
\end{tikzpicture}

\begin{center}
\vspace*{1.2cm}
{\fontsize{24}{30}\selectfont\color{white}\textbf{多种算法求解\\[0.25cm]0/1背包问题}}

\vspace{0.25cm}
\textcolor{ember}{\rule{2.8cm}{1.6pt}}

\vspace{0.6cm}
{\large\color{cold}性能比较实验}

\vspace{1.2cm}
{\normalsize\color{rule}贾静潇 \quad 张诗淇}

\vspace{0.25cm}
{\footnotesize\color{rule}算法分析与设计 · 2.3 问题求解类}

\vspace{0.4cm}
{\tiny\color{rule!60}数据集: FSU · Pi · JJ \quad 算法: 贪心 · DP · 回溯 · 分支限界}
\end{center}
\end{frame}
}

% ====== Slide 2: 问题定义 ======
\begin{frame}{0/1背包问题}

\begin{leftblock}{形式化定义}
给定 $n$ 个物品（重量 $w_i$、价值 $v_i$）和容量 $W$，
\[\text{max}\sum_{i\in S} v_i \qquad \text{s.t.}\;\sum_{i\in S} w_i \leq W,\; x_i \in \{0,1\}\]
\end{leftblock}

\vspace{0.2cm}

\begin{columns}[T]
\begin{column}{0.5\textwidth}
{\small\color{ink}\textbf{理论性质}}\\[0.1cm]
{\footnotesize\color{slate}NP-hard · 最优子结构 · 伪多项式可解}
\end{column}
\begin{column}{0.5\textwidth}
{\small\color{ink}\textbf{应用场景}}\\[0.1cm]
{\footnotesize\color{slate}资源分配 · 投资决策 · 货物装载 · 带宽分配}
\end{column}
\end{columns}

\vspace{0.3cm}
\begin{itemize}
  \item 每件物品\hi{要么整取、要么不取}，不可分割
  \item 物品间\hi{无先后顺序}约束，仅受容量限制
  \item 目标：在 $2^n$ 种组合中找到满足容量约束的最大价值子集
\end{itemize}

\end{frame}

% ====== Slide 3: 四种算法总览 ======
\begin{frame}{四种求解策略}

\vspace{-0.1cm}
\begin{table}
\renewcommand{\arraystretch}{1.25}
\small
\begin{tabular}{p{1.8cm}p{2.5cm}p{2.2cm}p{2.2cm}}
\toprule
\textbf{算法} & \textbf{核心思路} & \textbf{时间} & \textbf{最优？} \\
\midrule
\hi{贪心}      & 性价比排序，依次塞入   & $O(n\log n)$  & \hi{否} \\
\blu{动态规划} & 自底向上递推填表       & $O(nW)$       & 是 \\
回溯          & 深搜 + 上界剪枝       & \tb           & \tb \\
分支限界      & 优先队列 + 上下界     & \tb           & \tb \\
\bottomrule
\end{tabular}
\end{table}

\vspace{0.15cm}
\begin{center}
\begin{tikzpicture}[font=\footnotesize]
  \node[draw=rule,fill=soft,text width=8cm,align=center,rounded corners=3pt,inner sep=6pt] {
    {\color{ink}\textbf{同一问题 · 四种策略 · 不同瓶颈}}\\[0.08cm]
    {\color{slate}\small 贪心怕数据分布 \quad DP 怕大容量 \quad 搜索怕大物品数}
  };
\end{tikzpicture}
\end{center}

\end{frame}

% ====== Slide 4: 贪心 ======
\begin{frame}{贪心算法}

\begin{columns}[T]
\begin{column}{0.5\textwidth}

\begin{leftblock}{策略}
$r_i \gets v_i/w_i$，按 $r_i$ \hi{降序}排列，依次尝试放入——够就塞，跳过不回头。
\end{leftblock}

\vspace{0.1cm}
\begin{itemize}
  \item 时间 $O(n\log n)$ · 空间 $O(n)$
  \item 与 $W$ \hi{完全无关}
  \item 大 $n$ 下依然毫秒级
\end{itemize}

\end{column}
\begin{column}{0.5\textwidth}

\begin{leftblock}{失效示例}
容量 10
\vspace{0.05cm}
\begin{tabular}{cccc}
\hi{A} & w=6 & v=9  & $r$=1.5 \\
B     & w=5 & v=7  & $r$=1.4 \\
C     & w=5 & v=7  & $r$=1.4 \\
\end{tabular}
\par\vspace{0.1cm}
贪心拿 A：价值 9（剩余 4）\\
\blu{最优 B+C：价值 14}
\end{leftblock}

\end{column}
\end{columns}

\vspace{0.2cm}
\begin{center}
{\small\color{slate}每次选当前最优 → 局部最优不保证全局最优 → \hi{需要后悔，但不后悔}}
\end{center}

\end{frame}

% ====== Slide 5: DP ======
\begin{frame}{动态规划}

\begin{columns}[T]
\begin{column}{0.5\textwidth}
\begin{leftblock}{状态与转移}
\[\boxed{dp[w] = \max\{\, dp[w],\; dp[w-w_i] + v_i \,\}}\]
$dp[w]$ = 容量 $w$ 时的最大价值
\end{leftblock}

\vspace{0.1cm}
\begin{itemize}
  \item 内层 $w$ 必须\hi{倒序}
  \item 正序 → 物品可重复 → 完全背包
\end{itemize}
\end{column}

\begin{column}{0.5\textwidth}
\begin{leftblock}{复杂度}
\textbf{时间：} \hi{$O(nW)$} 伪多项式 \\
\textbf{空间：} $O(W)$（一维优化）

\vspace{0.12cm}
{\footnotesize\color{slate}
$n$=2000, $W$=250万：$\sim$50亿次迭代 → 10s \\
$W$=100亿：dp数组$\sim$40GB → \hi{OOM}
}
\end{leftblock}
\end{column}
\end{columns}

\vspace{0.2cm}
\begin{center}
{\small\color{slate}以 $W$ 为代价换最优解 · 当 $W$ 失控时从优势变为劣势}
\end{center}

\end{frame}

% ====== Slide 6: 回溯与分支限界 ======
\begin{frame}{回溯与分支限界}

\begin{columns}[T]
\begin{column}{0.5\textwidth}
\begin{leftblock}{回溯}
深度优先搜索解空间树

\vspace{0.08cm}
{\footnotesize\color{slate}
剪枝条件：剩余容量 + 价值上界 \\
最坏 $O(2^n)$，$n$ 大时不可行
}
\end{leftblock}
\end{column}
\begin{column}{0.5\textwidth}
\begin{leftblock}{分支限界}
优先队列驱动搜索

\vspace{0.08cm}
{\footnotesize\color{slate}
上界+下界联合剪枝 \\
优先扩展最有希望的节点
}
\end{leftblock}
\end{column}
\end{columns}

\vspace{0.3cm}
\begin{center}
{\large\color{slate}\tb}

\vspace{0.15cm}
{\tiny\color{rule}（策略描述、复杂度分析、实验结果待张诗淇补充）}
\end{center}

\end{frame}

% ====== Slide 7: 数据集 ======
\begin{frame}{三个数据源}

\vspace{-0.2cm}
\begin{table}
\renewcommand{\arraystretch}{1.2}
\footnotesize
\begin{tabular}{p{1.5cm}p{2.3cm}p{3.8cm}r}
\toprule
\textbf{数据源} & \textbf{来源} & \textbf{特征} & \textbf{数量} \\
\midrule
\hi{FSU}  & Florida State U. & 官方基准，$n$=5--24，附最优解对照 & 8 \\
\blu{Pi}  & 自生成(Pisinger) & 3种相关性 $\times$ 3种规模，固定种子 & 9 \\
\hi{JJ}  & JorikJooken (GH) & 学术标准集，$n$=400--1000，$W$至100亿 & 6 \\
\bottomrule
\end{tabular}
\end{table}

\vspace{0.1cm}
\begin{center}
\begin{tikzpicture}[font=\footnotesize]
  \node[draw=accent,fill=accent!8,text width=2.8cm,align=center,rounded corners=3pt,inner sep=4pt] (f) {\textbf{FSU}\\[0.03cm]\tiny 正确性验证};
  \node[draw=accent,fill=accent!8,text width=2.8cm,align=center,rounded corners=3pt,inner sep=4pt,right=0.3cm of f] (p) {\textbf{Pi}\\[0.03cm]\tiny 控制变量分析};
  \node[draw=accent,fill=accent!8,text width=2.8cm,align=center,rounded corners=3pt,inner sep=4pt,right=0.3cm of p] (j) {\textbf{JJ}\\[0.03cm]\tiny 大规模实战测试};
  \draw[->,ember,line width=1pt] (f) -- (p);
  \draw[->,ember,line width=1pt] (p) -- (j);
\end{tikzpicture}
\end{center}

\end{frame}

% ====== Slide 8: FSU 结果 ======
\begin{frame}{FSU 实验结果}

\vspace{-0.25cm}
\begin{center}
\begin{tikzpicture}
\begin{axis}[
  barplot, height=3.8cm, width=0.9\textwidth,
  ylabel={总价值},
  symbolic x coords={P01,P02,P03,P04,P05,P06,P07,P08},
  xtick=data,
  x tick label style={font=\fontsize{4}{4.5}\selectfont},
  legend style={font=\fontsize{4}{4.5}\selectfont,at={(0.98,0.98)},anchor=north east},
]
  \addplot[fill=ember!80,draw=ember] coordinates {
    (P01,309)(P02,47)(P03,146)(P04,102)(P05,858)(P06,1478)(P07,1441)(P08,13415886)
  };
  \addplot[fill=accent!80,draw=accent] coordinates {
    (P01,309)(P02,51)(P03,150)(P04,107)(P05,900)(P06,1735)(P07,1458)(P08,13549094)
  };
  \legend{贪心, DP/最优}
\end{axis}
\end{tikzpicture}
\end{center}

\vspace{-0.15cm}
\begin{itemize}
  \item \blu{DP 全 8 例命中最优} · 贪心仅 P01 达最优 · 最差 gap \hi{14.81\%}（P06）
  \item P08: $n$=24, $W$=640万，DP 1.59s，贪心 $<$0.2ms → \hi{万倍速差}
\end{itemize}

\end{frame}

% ====== Slide 9: Pi 结果 ======
\begin{frame}{Pi 实验结果}

\vspace{-0.25cm}
\begin{columns}[T]
\begin{column}{0.48\textwidth}
{\small\color{ink}\textbf{DP 耗时趋势（对数）}}\\[0.05cm]
\begin{tikzpicture}
\begin{axis}[
  narrowline, height=3.0cm,
  ymode=log, log basis y={10},
  xlabel={}, ylabel={耗时(ms)},
  xtick={0,1,2}, xticklabels={S(50),M(500),L(2000)},
  cycle list={{accent,mark=*},{ember,mark=*},{ink,mark=*}},
]
  \addplot coordinates {(0,10.18)(1,760.54)(2,9859.19)};
  \addlegendentry{不相关}
  \addplot coordinates {(0,0.99)(1,566.29)(2,9069.08)};
  \addlegendentry{弱相关}
  \addplot coordinates {(0,2.10)(1,527.93)(2,8244.92)};
  \addlegendentry{强相关}
\end{axis}
\end{tikzpicture}
\end{column}
\begin{column}{0.52\textwidth}
{\small\color{ink}\textbf{贪心 Gap 按模式}}\\[0.05cm]
\begin{tikzpicture}
\begin{axis}[
  narrowbar, height=3.0cm,
  ylabel={Gap (\%)},
  symbolic x coords={不相关,弱相关,强相关},
  xtick=data,
  x tick label style={font=\fontsize{4}{4.5}\selectfont},
]
  \addplot[fill=accent!60,draw=accent] coordinates {(不相关,0.00)(弱相关,0.00)(强相关,0.85)};
  \addlegendentry{S(50)}
  \addplot[fill=accent!85,draw=accent] coordinates {(不相关,0.01)(弱相关,0.04)(强相关,0.12)};
  \addlegendentry{M(500)}
  \addplot[fill=accent,draw=accent]  coordinates {(不相关,0.00)(弱相关,0.00)(强相关,0.02)};
  \addlegendentry{L(2000)}
\end{axis}
\end{tikzpicture}
\end{column}
\end{columns}

\vspace{-0.1cm}
{\footnotesize\color{slate}\hi{强相关 S 级 gap 最高 0.85\%} — 性价比一致时贪心排序失效 · DP S→L 耗时 $\times$1000 · 回溯/分支限界 \tb}

\end{frame}

% ====== Slide 10: JJ 结果 ======
\begin{frame}{JJ 实验结果}

\vspace{-0.25cm}
\begin{center}
{\small\color{ink}\textbf{DP vs 贪心耗时（对数坐标）}}
\begin{tikzpicture}
\begin{axis}[
  barplot, height=4.0cm, width=0.88\textwidth,
  ymode=log, log basis y={10},
  ylabel={耗时(ms)},
  symbolic x coords={n=400,f=0.1, n=400,f=0.3, n=1000,f=0.1, n=400,W=1亿},
  xtick=data,
  enlarge x limits=0.15,
  x tick label style={font=\fontsize{4}{4.5}\selectfont},
]
  \addplot[fill=ember!80,draw=ember] coordinates {
    (n=400,f=0.1, 51.67)(n=400,f=0.3, 0.50)(n=1000,f=0.1, 1.65)(n=400,W=1亿, 0.53)
  };
  \addlegendentry{贪心}
  \addplot[fill=accent!80,draw=accent] coordinates {
    (n=400,f=0.1, 1128)(n=400,f=0.3, 756)(n=1000,f=0.1, 1826)(n=400,W=1亿, 113019)
  };
  \addlegendentry{DP}
\end{axis}
\end{tikzpicture}
\end{center}

\vspace{-0.15cm}
\begin{itemize}
  \item $W$=1亿：DP \hi{113s} vs 贪心 $<$1ms，\textbf{11 万倍差距}
  \item $W$=100亿：DP \hi{OOM}（$\sim$40GB），贪心正常 · 回溯 $2^{400}$ → \hi{不可行}
\end{itemize}

\end{frame}

% ====== Slide 11: 核心洞察 ======
\begin{frame}{核心洞察：$n$ 与 $W$ 的博弈}

\vspace{-0.15cm}
\begin{center}
\renewcommand{\arraystretch}{2.0}
\large
\begin{tabular}{c|c|c|}
\multicolumn{1}{c}{} & \multicolumn{1}{c}{\textbf{$n$ 小}} & \multicolumn{1}{c}{\textbf{$n$ 大}} \\
\cline{2-3}
\textbf{$W$ 小} & \color{slate}都轻松 & \blu{DP 统治} \\
\cline{2-3}
\textbf{$W$ 大} & \hi{回溯/分支限界优} & \hi{贪心唯一可行} \\
\cline{2-3}
\end{tabular}
\end{center}

\vspace{0.25cm}

\begin{columns}[T]
\begin{column}{0.52\textwidth}
{\footnotesize\color{slate}
DP 的敌人是 \hi{$W$}：$W$ 翻 10 倍，耗时翻 10 倍。\\
回溯的敌人是 \hi{$n$}：$n$ 增 1，搜索空间翻倍。\\
贪心\blu{两者都不怕}——但以牺牲最优性为代价。
}
\end{column}
\begin{column}{0.44\textwidth}
\begin{center}
\textcolor{ember}{\rule{2cm}{0.5pt}}\\[0.2cm]
{\large\color{ink}\textbf{没有最优算法\\只有最优选择}}\\[0.2cm]
\textcolor{ember}{\rule{2cm}{0.5pt}}
\end{center}
\end{column}
\end{columns}

\end{frame}

% ====== Slide 12: 复杂度总表 ======
\begin{frame}{四种算法复杂度对比}

\vspace{-0.2cm}
\begin{table}
\renewcommand{\arraystretch}{1.3}
\small
\begin{tabular}{p{1.9cm}p{1.8cm}p{1.8cm}p{1.8cm}p{1.8cm}}
\toprule
\textbf{维度} & \textbf{贪心} & \textbf{DP} & \textbf{回溯} & \textbf{分支限界} \\
\midrule
时间复杂度 & $O(n\log n)$ & \hi{$O(nW)$} & \tb & \tb \\
空间复杂度 & $O(n)$ & \hi{$O(W)$} & \tb & \tb \\
保证最优？ & \hi{否} & 是 & \tb & \tb \\
敏感于 & 数据分布 & $W$ & $n$ & \tb \\
大容量可行？ & \blu{始终} & \hi{否($>10^8$)} & \tb & \tb \\
\bottomrule
\end{tabular}
\end{table}

\vspace{0.15cm}
\begin{center}
{\large\color{ink}贪心：快但不准 \quad$\mid$\quad DP：准但怕大W \quad$\mid$\quad 搜索类：\tb}\\[0.2cm]
{\footnotesize\color{slate}四种算法构成连续谱系——从极速近似到精确最优，选型取决于 $n$、$W$ 和精度约束。}
\end{center}

\end{frame}

% ====== Slide 13: 总结 ======
\begin{frame}{总结}

\vspace{-0.1cm}
\begin{enumerate}
  \item 0/1背包问题以\hi{简洁的结构承载了丰富的算法策略差异}

  \item \hi{贪心} — 毫秒级 · 与 $W$ 无关 · gap 最高 15\% · \blu{实时场景首选}

  \item \blu{DP} — 始终最优 · $O(nW)$ 伪多项式 · $W>10^8$ 不可用 · \blu{精度优先}

  \item 回溯/分支限界 — $n$ 小 $W$ 大时优于 DP · 剪枝策略决定效率 \tb

  \item \textbf{选型公式} — 看 $n$、看 $W$、看精度需求 → \hi{三种瓶颈各不同}
\end{enumerate}

\vspace{0.3cm}
\begin{center}
{\large\color{ink}\textbf{时间与最优性的折中，是算法设计的永恒主题。}}
\end{center}

\vspace{0.15cm}
\begin{flushright}
{\footnotesize\color{rule}谢谢！}
\end{flushright}

\end{frame}

\end{document}
